\chapter{Lesson 6: Making Plans}

This lesson moves into the future: saying what you are going to do, suggesting things
to do together, and agreeing on a time and place. The grammar is the future/intention
ending \ko{\til (으)ㄹ 거예요}, the suggestion question \ko{\til (으)ㄹ까요?}, the
``let's'' form \ko{\til (으)ㅂ시다}, and the promise ending \ko{\til (으)ㄹ게요}.

\section{Vocabulary}

\begin{table}[H]
\centering
\begin{tabular}{lll}
\toprule
Korean & RR & English \\
\midrule
\ko{계획} & gyehoek & plan \\
\ko{약속} & yaksok & appointment, promise, plans with someone \\
\ko{방학} & banghak & school vacation \\
\ko{여행하다 / 여행 가다} & yeohaenghada / yeohaeng gada & to travel / go on a trip \\
\ko{초대하다} & chodaehada & to invite \\
\ko{파티} & pati & party \\
\ko{이번 주말} & ibeon jumal & this weekend \\
\ko{다음 달} & daeum dal & next month \\
\ko{내년} & naenyeon & next year \\
\ko{나중에} & najunge & later \\
\ko{언제} & eonje & when \\
\ko{어디서} & eodiseo & where (short for \ko{어디에서}) \\
\ko{시간(이) 있다} & sigan(i) itda & to be free, have time \\
\ko{바쁘다} & bappeuda & to be busy \\
\ko{괜찮다} & gwaenchanta & to be OK, fine \\
\ko{좋아요} & joayo & sounds good, OK \\
\ko{글쎄요} & geulsseyo & well\ldots, I'm not sure \\
\bottomrule
\end{tabular}
\end{table}

\section{Grammar 1: \ko{\til (으)ㄹ 거예요} (going to)}

\ko{\til (으)ㄹ 거예요} expresses a \textbf{plan or intention} (with \ko{저/우리} or in
questions to \ko{you}) and a \textbf{prediction} (with third-person subjects, weather,
etc.).

\begin{itemize}
  \item Vowel stem: \ko{ㄹ 거예요}. \ko{가다} $\to$ \ko{갈 거예요} gal geoyeyo.
  \item Consonant stem: \ko{을 거예요}. \ko{먹다} $\to$ \ko{먹을 거예요} meogeul
        geoyeyo.
  \item \ko{ㄹ} stem: \ko{거예요} only. \ko{놀다} $\to$ \ko{놀 거예요} nol geoyeyo.
  \item Irregulars: \ko{듣다} $\to$ \ko{들을 거예요}; \ko{덥다} $\to$ \ko{더울 거예요}.
\end{itemize}

\begin{remark}
The \ko{거} after the \ko{ㄹ} modifier is tensed in speech: \ko{할 거예요} is
[\ko{할 꺼에요}] (RR still writes hal geoyeyo), since \ko{예} after a consonant
is pronounced [\ko{에}]. Write it \ko{거예요}, never \ko{꺼예요} or \ko{거에요}.
\end{remark}

\begin{center}
\begin{tikzpicture}[x=1cm, every node/.style={font=\small}]
  \draw[-{Stealth}, thick] (-0.5,0) -- (11,0);
  \fill[red!60] (5,0) circle (3pt);
  \node[above, font=\scriptsize, red!60!black] at (5,0.15) {now};
  \node[draw, rounded corners, fill=gray!10, align=center] at (1.6,1.1) {\ko{봤어요}\\ \scriptsize watched};
  \node[draw, rounded corners, fill=blue!8, align=center] at (5,1.6) {\ko{봐요}\\ \scriptsize watch (habit / now)};
  \node[draw, rounded corners, fill=orange!15, align=center] at (5,-1.2) {\ko{보고 있어요}\\ \scriptsize am watching (Lesson 7)};
  \node[draw, rounded corners, fill=green!12, align=center] at (8.6,1.1) {\ko{볼 거예요}\\ \scriptsize will / going to watch};
  \node[font=\scriptsize] at (1.6,-0.35) {past};
  \node[font=\scriptsize] at (8.6,-0.35) {future};
  \draw[dashed] (1.6,0.65) -- (1.6,0);
  \draw[dashed] (8.6,0.65) -- (8.6,0);
  \draw[dashed] (5,1.15) -- (5,0.1);
  \draw[dashed] (5,-0.75) -- (5,-0.1);
\end{tikzpicture}
\end{center}

\begin{example}[Plans and predictions]
Plans (first person) and predictions (weather, other people):
\begin{itemize}
  \item \ko{이번 주말에 뭐 할 거예요?} Ibeon jumare mwo hal geoyeyo? (What are you going
        to do this weekend?)
  \item \ko{토요일에 친구 생일 파티에 갈 거예요.} Toyoire chingu saengil patie gal
        geoyeyo. (On Saturday I'm going to a friend's birthday party.)
  \item \ko{방학에 한국에 여행 갈 거예요. 부산에서 바다를 볼 거예요.} Banghage hanguge
        yeohaeng gal geoyeyo. Busaneseo badareul bol geoyeyo. (Over the break I'm going to
        travel to Korea. I'll see the ocean in Busan.)
  \item \ko{내일은 비가 올 거예요.} Naeireun biga ol geoyeyo. (It'll probably rain
        tomorrow.)
  \item \ko{민수 씨는 지금 바쁠 거예요.} Minsu ssineun jigeum bappeul geoyeyo. (Minsu is
        probably busy now.)
  \item \ko{할머니께서 내일 오실 거예요.} Halmeonikkeseo naeil osil geoyeyo. (Grandma
        will come tomorrow.) --- honorific \ko{시} + \ko{ㄹ 거예요}.
\end{itemize}
\end{example}

\section{Grammar 2: \ko{\til (으)ㄹ까요?} (shall we / shall I?)}

\ko{\til (으)ㄹ까요?} attaches to the stem the same way as \ko{\til (으)ㄹ 거예요}. It has
three uses:
\begin{enumerate}
  \item \textbf{Suggestion} with \ko{우리} (shall we?): \ko{우리 같이 점심 먹을까요?}
        Uri gachi jeomsim meogeulkkayo? (Shall we have lunch together?)
  \item \textbf{Offer} with \ko{제가} (shall I?): \ko{제가 창문을 열까요?} Jega
        changmuneul yeolkkayo? (Shall I open the window?)
  \item \textbf{Wondering} about a third person or event (do you think \ldots?):
        \ko{내일 날씨가 좋을까요?} Naeil nalssiga joeulkkayo? (Do you think the weather
        will be nice tomorrow?)
\end{enumerate}

Typical replies to a suggestion: \ko{좋아요!} (sounds good), \ko{그래요} (let's),
\ko{미안해요. 그날은 약속이 있어요.} (sorry, I have plans that day), or the noncommittal
\ko{글쎄요\ldots}.

\section{Grammar 3: \ko{\til (으)ㅂ시다} (let's)}

\ko{\til (으)ㅂ시다} is the ``let's'' form. Vowel stem: \ko{ㅂ시다}; consonant stem:
\ko{읍시다}; \ko{ㄹ} stem: drop \ko{ㄹ}, add \ko{ㅂ시다}.

\begin{table}[H]
\centering
\begin{tabular}{llll}
\toprule
Dictionary & \ko{\til (으)ㄹ까요?} & \ko{\til (으)ㅂ시다} & Meaning \\
\midrule
\ko{가다} & \ko{갈까요?} & \ko{갑시다} & go \\
\ko{먹다} & \ko{먹을까요?} & \ko{먹읍시다} & eat \\
\ko{만나다} & \ko{만날까요?} & \ko{만납시다} & meet \\
\ko{만들다} & \ko{만들까요?} & \ko{만듭시다} & make \\
\ko{걷다} & \ko{걸을까요?} & \ko{걸읍시다} & walk (\ko{ㄷ}-irregular) \\
\bottomrule
\end{tabular}
\end{table}

\begin{important}[Do not say \ko{\til (으)ㅂ시다} to a superior]
Although it ends in a formal-looking \ko{\til 다}, \ko{\til (으)ㅂ시다} sounds like the
speaker is leading the group. Use it with peers and juniors, or in announcements
(\ko{조용히 합시다}, ``let's keep quiet''). To suggest something to a teacher or an older
person, use \ko{\til (으)실까요?} (\ko{같이 가실까요?}) or simply the \ko{해요} form
with rising-then-falling intonation: \ko{같이 가요.} (Let's go together.) In \ko{해요}체
the statement, question, request and suggestion are all the same form; only intonation
and context tell them apart.
\end{important}

\section{Grammar 4: \ko{\til (으)ㄹ게요} (I'll do it --- a promise)}

\ko{\til (으)ㄹ게요} is used only with a first-person subject, to commit to doing something
\emph{in response to} the listener (an offer, promise or decision that takes them into
account). Compare:
\begin{itemize}
  \item \ko{제가 커피 살게요.} Jega keopi salgeyo. (I'll buy the coffee --- I'm offering.)
  \item \ko{내일 커피를 살 거예요.} Naeil keopireul sal geoyeyo. (I'm going to buy coffee
        tomorrow --- just reporting my plan.)
  \item \ko{이따가 전화할게요.} Ittaga jeonhwahalgeyo. (I'll call you later.)
\end{itemize}

\section{Agreeing on a time and place}

\begin{example}[Setting up plans]
A typical exchange between two classmates:
\begin{quote}
\ko{지수: 수진 씨, 이번 토요일에 시간 있어요?}\\
\ko{수진: 네, 있어요. 왜요?}\\
\ko{지수: 같이 영화 볼까요? 새 영화가 재미있을 거예요.}\\
\ko{수진: 좋아요! 몇 시에 만날까요?}\\
\ko{지수: 두 시 어때요? 영화관 앞에서 만나요.}\\
\ko{수진: 그래요. 영화 보고 저녁도 먹읍시다. 제가 살게요.}
\end{quote}
Jisu: Sujin, are you free this Saturday? / Sujin: Yes, why? / Jisu: Shall we see a movie
together? The new movie should be good. / Sujin: Great! What time shall we meet? / Jisu:
How about two o'clock? Let's meet in front of the theatre. / Sujin: OK. Let's watch the
movie and eat dinner too. I'll pay.
\end{example}

\begin{center}
\begin{tikzpicture}[every node/.style={font=\scriptsize}]
  \foreach \i/\d in {0/\ko{월}, 1/\ko{화}, 2/\ko{수}, 3/\ko{목}, 4/\ko{금}, 5/\ko{토}, 6/\ko{일}} {
    \node[draw, fill=gray!15, minimum width=1.9cm, minimum height=0.55cm] at (\i*1.9,0) {\d};
    \node[draw, minimum width=1.9cm, minimum height=1.3cm] at (\i*1.9,-0.93) {};
  }
  \node[align=center] at (0,-0.93) {\ko{수업}\\ \ko{아르바이트}};
  \node[align=center] at (1.9,-0.93) {\ko{수업}};
  \node[align=center] at (3.8,-0.93) {\ko{시험!}};
  \node[align=center] at (5.7,-0.93) {\ko{수업}\\ \ko{동아리}};
  \node[align=center] at (7.6,-0.93) {\ko{친구 생일}\\ \ko{파티}};
  \node[align=center, fill=yellow!40] at (9.5,-0.93) {\ko{2시 영화}\\ \ko{지수 씨}};
  \node[align=center] at (11.4,-0.93) {\ko{청소, 쉬기}};
  \node[font=\small] at (5.7,0.65) {\ko{이번 주 계획} (this week's plans)};
\end{tikzpicture}
\end{center}

\begin{example}[Describing the week]
Reading the planner above aloud:
\begin{itemize}
  \item \ko{수요일에 시험이 있어서 화요일에 도서관에서 공부할 거예요.} Suyoire siheomi
        isseoseo hwayoire doseogwaneseo gongbuhal geoyeyo. (I have an exam on Wednesday,
        so on Tuesday I'll study at the library.)
  \item \ko{토요일 두 시에 지수 씨하고 영화를 볼 거예요.} Toyoil du sie Jisu ssihago
        yeonghwareul bol geoyeyo. (At two on Saturday I'm going to see a movie with
        Jisu.)
  \item \ko{일요일에는 집을 청소하고 쉴 거예요.} Iryoireneun jibeul cheongsohago swil
        geoyeyo. (On Sunday I'll clean the house and rest.)
\end{itemize}
\end{example}

\begin{aside}
\textbf{``Let's eat sometime.''} The phrase \ko{언제 밥 한번 먹어요} (let's grab a meal
sometime) is often a friendly way of ending a conversation rather than a real
invitation, much like ``we should get coffee sometime'' in English. A real plan has a
date and time. When friends go out, the bill is often paid by one person (frequently the
oldest, or the \ko{선배}), with the others paying for the next round of coffee or drinks
(\ko{2차} icha, the ``second round''), though splitting the bill (\ko{더치페이}, from
``Dutch pay'') has become common among younger people.
\end{aside}

\section*{Review summary}
\begin{itemize}
  \item \ko{\til (으)ㄹ 거예요}: plan (I/we) or prediction (others, weather).
  \item \ko{\til (으)ㄹ까요?}: shall we? (\ko{우리}), shall I? (\ko{제가}), do you think?
        (others).
  \item \ko{\til (으)ㅂ시다}: let's --- only to peers and juniors; to seniors use
        \ko{\til 아/어요} or \ko{\til (으)실까요?}.
  \item \ko{\til (으)ㄹ게요}: I'll do it (promise/offer, responsive to the listener).
  \item \ko{ㄹ}-stems: \ko{놀 거예요, 만들까요?, 만듭시다}.
\end{itemize}
